% from ega3-1-fr.tex
\oldpage[III]{108}
(0, 12.2.2), et comme (\hyperref[III.2.4.1.2-fr]{2.4.1.2}) provient
lui-même (par application de (0\textsubscript{I}, 4.2.2.1)) de
l'homomorphisme
\[
\mathcal O_{X'}\otimes\widetilde{\mathcal M}
\longrightarrow\widetilde{\mathcal M}
\]
définissant la structure de $\mathcal O_{X'}$-Module de
$\widetilde{\mathcal M}$, le diagramme
\[
\xymatrix{
f_*(g'_*(\mathcal O_{X'}))\otimes
R^pf_*(g'_*(\widetilde{\mathcal M}))
\ar[r]\ar[d] &
R^pf_*(g'_*(\widetilde{\mathcal M}))\ar[d]\\
h_*(\mathcal O_{X'})\otimes R^ph_*(\widetilde{\mathcal M})
\ar[r] & R^ph_*(\widetilde{\mathcal M})
}
\]
est commutatif (0, 12.2.6)~; composant les flèches horizontales avec
l'homomorphisme provenant de l'homomorphisme canonique
\[
\mathcal S\longrightarrow f_*(f^*(\mathcal S))=f_*(\mathcal S')
=f_*(g'_*(\mathcal O_{X'}))=h_*(\mathcal O_{X'}),
\]
on obtient notre assertion. D'autre part, puisque $g$ est affine et que
$f'$ est séparé et quasi-compact, l'homomorphisme canonique
$R^ph_*(\widetilde{\mathcal M})\to
g_*(R^pf'_*(\widetilde{\mathcal M}))$ est bijectif
(\hyperref[III.1.4.14-fr]{1.4.14}), et on démontre comme ci-dessus que
c'est un isomorphisme de $\mathcal S$-Modules (en utilisant cette fois la
commutativité de (0, 12.2.6.2)). Or, $f'$ étant projectif et
$\widetilde{\mathcal M}$ cohérent,
$R^pf'_*(\widetilde{\mathcal M})$ est un
$\mathcal O_{Y'}$-Module cohérent en vertu de
(\hyperref[III.2.2.1-fr]{2.2.1})~; on en conclut que
$g_*(R^pf'_*(\widetilde{\mathcal M}))$ est un $\mathcal S$-Module de type
fini (II, 1.4.5).

(ii) Soit $\mathcal L'=g'^*(\mathcal L)$, qui est un
$\mathcal O_{X'}$-Module inversible~; pour tout $n\in\mathbb Z$, on a
\[
g'_*(\widetilde{\mathcal M}\otimes\mathcal L'^{\otimes n})
=g'_*(\widetilde{\mathcal M})\otimes\mathcal L^{\otimes n}
=\mathcal M(n)
\]
(0\textsubscript{I}, 5.4.10) à un isomorphisme près~; on peut appliquer
à $\widetilde{\mathcal M}\otimes\mathcal L'^{\otimes n}$ le
raisonnement fait dans (i) pour $\widetilde{\mathcal M}$, qui prouve que
\[
R^pf_*(g'_*(\widetilde{\mathcal M}\otimes\mathcal L'^{\otimes n}))
\quad\hbox{est isomorphe à}\quad
g_*(R^pf'_*(\widetilde{\mathcal M}\otimes\mathcal L'^{\otimes n})).
\]
Or $\mathcal L'$ est ample pour $f'$ (II, 4.6.13, (iii)), donc il résulte
de (\hyperref[III.2.2.1-fr]{2.2.1}) qu'il existe un entier $N$ tel que
\[
R^pf'_*(\widetilde{\mathcal M}\otimes\mathcal L'^{\otimes n})=0
\]
pour tout $p$ et tout $n\geq N$, ce qui prouve
(\hyperref[III.2.4.1.1-fr]{2.4.1.1}). Enfin, il résulte encore de
(\hyperref[III.2.2.1-fr]{2.2.1}) qu'on peut supposer $N$ tel que pour
$n\geq N$, l'homomorphisme canonique
\[
f'^*(f'_*(\widetilde{\mathcal M}\otimes\mathcal L'^{\otimes n}))
\longrightarrow\widetilde{\mathcal M}\otimes\mathcal L'^{\otimes n}
\]
soit surjectif~; comme $g'_*$ est un foncteur exact (II, 1.4.4),
l'homomorphisme correspondant
\[
g'_*(f'^*(f'_*(\widetilde{\mathcal M}\otimes
\mathcal L'^{\otimes n})))
\longrightarrow
g'_*(\widetilde{\mathcal M}\otimes\mathcal L'^{\otimes n})
=\mathcal M(n)
\]
est surjectif. Or, on a
\[
g'_*(f'^*(f'_*(\widetilde{\mathcal M}\otimes
\mathcal L'^{\otimes n})))
=f^*(g_*(f'_*(\widetilde{\mathcal M}\otimes
\mathcal L'^{\otimes n})))
\]
(II, 1.5.2) et comme $g_*\circ f'_* = f_*\circ g'_*$, on voit
finalement que l'on a
\[
g'_*(f'^*(f'_*(\widetilde{\mathcal M}\otimes\mathcal L'^{\otimes n})))
=f^*(f_*(g'_*(\widetilde{\mathcal M}\otimes
\mathcal L'^{\otimes n})))=f^*(f_*(\mathcal M(n))),
\]
ce qui achève la démonstration.
\end{proof}

\begin{env}[2.4.2]
\label{III.2.4.2-fr}
Nous aurons en particulier à appliquer
(\hyperref[III.2.4.1-fr]{2.4.1}) lorsque $\mathcal S$ est une
$\mathcal O_Y$-Algèbre graduée à degrés positifs,
$\mathcal M=\bigoplus_{k\in\mathbb Z}\mathcal M_k$ un
$\mathcal S'$-Module gradué. Alors (avec les

